\documentclass[11pt]{article}
\usepackage[T1]{fontenc}
\usepackage{lmodern}
\usepackage[utf8]{inputenc}
\usepackage{amsmath,amssymb,amsthm}
\usepackage[a4paper,margin=28mm]{geometry}
\usepackage{booktabs}
\usepackage[hidelinks]{hyperref}
\hypersetup{
  pdftitle={A Degree--Independent-Domination Bound for the Saturation Number: A Conjecture},
  pdfauthor={Hassine Achour},
  pdfsubject={Graph theory conjecture on saturation and independent domination}
}

\newtheorem{conjecture}{Conjecture}
\newtheorem{remark}{Remark}
\newcommand{\sat}{s}

\title{A Degree--Independent-Domination Bound for the Saturation Number: A Conjecture}
\author{Hassine Achour\\\texttt{hassine.achour@gmail.com}\\
\small\url{https://doi.org/10.5281/zenodo.21801900}}
\date{August 5, 2026}

\begin{document}
\maketitle

\begin{abstract}
For a finite simple graph $G$, let $\Delta(G)$ denote its maximum degree, let $i(G)$ be its independent domination number, and let $\sat(G)$ be its saturation number, namely the minimum cardinality of a maximal matching. We propose that every connected nontrivial graph satisfies
\[
  \sat(G)\leq \left\lceil\frac{\Delta(G)i(G)}{2}\right\rceil.
\]
The bound is sharp for complete graphs, while connectedness cannot be omitted. We report an exhaustive verification for all $273{,}192$ connected nonisomorphic graphs of orders $2$ through $9$, as well as supplementary random and regular-graph tests. These finite computations are evidence only and do not constitute a proof.
\end{abstract}

\section{Definitions and conjecture}
All graphs considered here are finite, simple, and undirected. A vertex set is \emph{independent dominating} if it is both independent and dominating. Its minimum possible cardinality is the \emph{independent domination number} $i(G)$ \cite{goddard2013}. A matching is \emph{maximal} if no edge can be added while preserving the matching property. The minimum cardinality of a maximal matching is the \emph{saturation number} $\sat(G)$, also denoted in the literature by parameters associated with minimum edge domination \cite{yannakakis1980,tavakoli2023}. Equivalently, $\sat(G)=i(L(G))=\gamma(L(G))$, where $L(G)$ is the line graph \cite[Theorem~4.1]{tavakoli2023}.

\begin{conjecture}\label{conj:main}
Every finite, simple, connected, nontrivial graph $G$ satisfies
\begin{equation}\label{eq:main}
 \boxed{\quad \sat(G)\leq \left\lceil\frac{\Delta(G)i(G)}{2}\right\rceil.\quad}
\end{equation}
\end{conjecture}

The conjecture connects a vertex parameter with an edge parameter. Existing work on minimum maximal matchings commonly gives bounds in terms of order, size, maximum degree, or restricted graph classes; see, for example, \cite{baste2019,cames2020}. To the best of the author's knowledge, no previously published bound relating $\sat(G)$ directly to $\Delta(G)$ and $i(G)$ in the form \eqref{eq:main} was identified in the literature searches performed through August 4, 2026. This is a limited priority statement, not a guarantee that an equivalent formulation does not occur in unindexed or differently notated literature.

\section{Sharpness and the role of connectedness}
The proposed bound is sharp for every positive value of the maximum degree. Indeed, for $G=K_{\Delta+1}$,
\[
 i(G)=1,\qquad
 \sat(G)=\left\lfloor\frac{\Delta+1}{2}\right\rfloor
       =\left\lceil\frac{\Delta}{2}\right\rceil.
\]
Thus equality holds in \eqref{eq:main}.

Connectedness is necessary for the statement as written. For the disjoint union $G=2K_2$,
\[
 \Delta(G)=1,\qquad i(G)=2,\qquad \sat(G)=2,
\]
whereas $\lceil\Delta(G)i(G)/2\rceil=1$. A componentwise variant would naturally require summing separate ceilings.

\section{Computational evidence}
The invariant $i(G)$ was computed exactly by branch-and-bound over independent dominating sets. The saturation number was computed through the identity $\sat(G)=i(L(G))$, with witness validation in the accompanying code. Isomorphism-free connected graphs were generated with \textsc{nauty} $2.8.8$ \cite{mckay2014}. The complete connected-graph campaign covered every order separately from $2$ through $9$.

\begin{center}
\begin{tabular}{lrrr}
\toprule
Campaign & Evaluations & Violations & Equalities \\
\midrule
All connected graphs, orders $2$--$9$ & $273{,}192$ & $0$ & $550$ \\
Random connected graphs, orders $10$--$16$ & $630$ & $0$ & $0$ \\
Regular-graph reports with pair data & $1{,}496{,}590$ & $0$ & $3{,}831$ \\
\bottomrule
\end{tabular}
\end{center}

The random campaign used Erd\H{o}s--R\'enyi graphs with probabilities $0.15$, $0.30$, and $0.50$, thirty connected samples per order--probability pair, and seed $20260804$. The regular-graph figure aggregates previously generated exhaustive, random, lift, and targeted reports containing exact $(i,\sat)$ pair distributions. These campaigns overlap and therefore must not be added to obtain a count of distinct graphs.

For every tested graph, the recorded slack was
\[
 \left\lceil\frac{\Delta(G)i(G)}2\right\rceil-\sat(G).
\]
The exhaustive connected campaign produced no negative slack. Reproducible scripts and summary outputs accompany this note. Because the general computation uses the line-graph identity, the reported search should be viewed as exact computational evidence rather than as a formally certified proof. In particular, no finite search establishes Conjecture~\ref{conj:main} for arbitrary order.

\section{Questions}
The immediate problem is to prove or disprove Conjecture~\ref{conj:main}. Natural subsidiary questions include characterizing equality, proving the bound for trees or bounded-degree classes, and identifying structural restrictions on a smallest counterexample.

\section*{Data and code availability}
The source, verification scripts, and machine-readable summaries are available from
\href{https://github.com/hermespromox/degree-independent-domination-saturation-conjecture}{GitHub}.

The archived preprint is available from
\href{https://doi.org/10.5281/zenodo.21801900}{Zenodo} (DOI: 10.5281/zenodo.21801900).

\begin{thebibliography}{9}
\bibitem{baste2019}
M. Baste, M. F\"urst, M. A. Henning, L. Mohr, and D. Rautenbach,
\emph{Bounding and approximating minimum maximal matchings in regular graphs},
arXiv:1905.12241, 2019.

\bibitem{cames2020}
W. Cames van Batenburg,
\emph{Minimum maximal matchings in cubic graphs},
arXiv:2008.01863, 2020.

\bibitem{goddard2013}
W. Goddard and M. A. Henning,
\emph{Independent domination in graphs: A survey and recent results},
Discrete Mathematics 313 (2013), 839--854.
\url{https://doi.org/10.1016/j.disc.2012.11.031}.

\bibitem{mckay2014}
B. D. McKay and A. Piperno,
\emph{Practical graph isomorphism, II},
Journal of Symbolic Computation 60 (2014), 94--112.
\url{https://doi.org/10.1016/j.jsc.2013.09.003}.

\bibitem{tavakoli2023}
M. Tavakoli and T. Do\v{s}li\'c,
\emph{Smallest maximal matchings of graphs},
Hacettepe Journal of Mathematics and Statistics 52(2) (2023), 356--366.
\url{https://doi.org/10.15672/hujms.1095437}.

\bibitem{yannakakis1980}
M. Yannakakis and F. Gavril,
\emph{Edge dominating sets in graphs},
SIAM Journal on Applied Mathematics 38(3) (1980), 364--372.
\url{https://doi.org/10.1137/0138030}.
\end{thebibliography}
\end{document}
